\begin{table}[htbp]
\centering\small
\caption{Lower bound on friction identified from each Argonne D3 drive cycle, taking the larger of the driven-axle traction utilisation and the all-wheel braking utilisation, and in brackets its fraction of the median friction $0.97$ pinned by the FMVSS~135 stops. The Focus Electric has no highway run.}\label{tab:d3cycles}
\begin{tabular}{lccc}
\toprule
Vehicle & UDDS (urban) & Highway & US06 (aggressive)\\
\midrule
2010 Toyota Prius & $0.32$ ($0.33$) & $0.32$ ($0.33$) & $0.72$ ($0.75$)\\
2012 Ford Focus & $0.36$ ($0.38$) & $0.35$ ($0.36$) & $0.69$ ($0.71$)\\
2012 Nissan Leaf & $0.36$ ($0.38$) & $0.29$ ($0.30$) & $0.60$ ($0.62$)\\
2013 Ford Focus Electric & $0.63$ ($0.65$) & n/a & $0.62$ ($0.64$)\\
\midrule
Median fraction & $0.38$ & $0.33$ & $0.68$\\
\bottomrule
\end{tabular}
\end{table}

\begin{table}[htbp]
\centering\small
\caption{Check of Theorem~\ref{thm:chance} on the rural route with the friction ellipse. For each vehicle and safety level $\delta$, the plan is the minimum-time profile at the friction multiplier $q(\delta)$ of the $114$ measured values, with trip time $T^\star$ (s). The plan was executed against all $114$ values. The table gives the number of values below $q(\delta)$, the number at which the plan was not executable, and the infeasible share, which is at most $\delta$ in every case.}\label{tab:app-chance}
\begin{tabular}{lcccccc}
\toprule
Vehicle & $\delta$ & $q(\delta)$ & $T^\star$ & Below & Infeasible & Share\\
\midrule
Leaf & $0.05$ & $0.805$ & $273.23$ & $5$ & $5$ & $0.044$\\
Leaf & $0.25$ & $0.940$ & $272.12$ & $28$ & $28$ & $0.246$\\
Civic & $0.05$ & $0.805$ & $273.90$ & $5$ & $5$ & $0.044$\\
Civic & $0.25$ & $0.940$ & $272.09$ & $28$ & $28$ & $0.246$\\
ZR1 & $0.05$ & $0.805$ & $269.46$ & $5$ & $5$ & $0.044$\\
ZR1 & $0.25$ & $0.940$ & $268.81$ & $28$ & $28$ & $0.246$\\
\bottomrule
\end{tabular}
\end{table}

\begin{table}[htbp]
\centering\small
\caption{Check of Proposition~\ref{prop:excess} on the urban route. For each vehicle and planning fraction $r$, the price (s) predicted by the transition-excess law, summed over the stop--go cycles, against the price computed by the three-pass solver, with their relative difference and the share of the predicted price due to braking.}\label{tab:app-excess}
\begin{tabular}{lcccccc}
\toprule
Vehicle & $r$ & Cycles & Law & Solver & Difference & Braking\\
\midrule
Leaf & $0.38$ & $13$ & $45.79$ & $45.77$ & $+0.0\%$ & $35\%$\\
Leaf & $0.68$ & $13$ & $11.11$ & $11.11$ & $+0.0\%$ & $42\%$\\
Civic & $0.38$ & $13$ & $54.97$ & $54.94$ & $+0.1\%$ & $31\%$\\
Civic & $0.68$ & $13$ & $15.70$ & $15.69$ & $+0.1\%$ & $32\%$\\
ZR1 & $0.38$ & $13$ & $36.21$ & $36.15$ & $+0.1\%$ & $30\%$\\
ZR1 & $0.68$ & $13$ & $10.32$ & $10.29$ & $+0.3\%$ & $30\%$\\
\bottomrule
\end{tabular}
\end{table}
